% ============================================================
% paper/sections/01_problem_failure_modes.tex
% Problem Statement and Structural Failure Modes
% ============================================================

\section{Problem Statement and Structural Failure Modes}

\subsection{Limits of Conventional Enterprise Analysis}

Enterprise analysis is commonly conducted under implicit assumptions
that remain largely unexamined. Among the most prevalent are the
assumptions that the enterprise persists as a coherent object, that
diagnosis remains meaningful under arbitrary degradation, and that
recovery is always, in principle, attainable.

These assumptions are rarely stated explicitly, yet they shape the
interpretation of observations and the continuation of analysis even
when structural coherence is no longer present. As a result, analytical
output may be produced in regimes where its validity is undefined.

Within such approaches, failure is typically framed as a local defect:
a gap in governance, an execution issue, or a misalignment between
strategy and operations. Structural conditions for existence are seldom
formalized.

\subsection{Enterprises as Conditional Objects}

In the Ontology-of-Continua perspective, an enterprise is not assumed to
exist unconditionally. It is treated as a continuum whose existence,
stability, and diagnosability depend on explicit structural conditions.

The enterprise continuum $K_{\mathrm{EA}}$ is admissible for diagnosis
only while:
\begin{itemize}
  \item its state space $\Omega_{\mathrm{EA}}$ is non-empty,
  \item its boundary $\partial\Omega_{\mathrm{EA}}$ remains well-defined,
  \item required cycles $C_{\mathrm{EA}}$ are present,
  \item observability and connectivity are sufficient to evaluate
        thresholds.
\end{itemize}

When these conditions are violated, continued analysis is no longer
defined within the model.

\subsection{Structural Failure Versus Performance Degradation}

A critical distinction made in $K_{\mathrm{EA}}$ is between performance
degradation and structural failure.

Performance degradation refers to states in which the enterprise
continues to satisfy existence conditions, but exhibits reduced
capability along one or more axes. Structural failure, by contrast,
refers to violations of the conditions that allow the enterprise to be
treated as a continuum at all.

Only the former admits continued diagnosis within $K_{\mathrm{EA}}$.

\subsection{Classes of Structural Failure Modes}

Structural failure modes are defined through threshold components
$f_k$ as specified in \texttt{ETS.K\_EA.v1.1}. These components do not
measure outcomes; they express violations of admissibility conditions.

The principal classes are:

\paragraph{Existence Failure ($f_{\text{exist}}$).}
The enterprise no longer maintains a non-empty admissible state space.
This may occur due to legal dissolution, loss of identity, or complete
loss of operational closure.

\paragraph{Stability Failure ($f_{\text{stab}}$).}
Recurrent cycles required for persistence cease to function. The system
may still exist momentarily, but cannot maintain itself over the
internal time scale $\tau_{\mathrm{EA}}$.

\paragraph{Inertial Failure ($f_{\text{inertia}}$).}
The enterprise remains structurally intact but is unable to transition
between admissible states without external intervention. Diagnosis
remains possible, but evolution is constrained.

\paragraph{Collapse Failure ($f_{\text{collapse}}$).}
Structural tension exceeds admissible bounds, leading to irreversible
loss of coherence. Some cycles or axes may persist locally, but the
enterprise as a whole no longer satisfies stability conditions.

\paragraph{Termination Condition ($f_{\text{stop}}$).}
Conditions under which diagnosis itself must terminate, regardless of
other indicators. This includes loss of observability and violation of
parent compatibility constraints.

\subsection{STOP as a Diagnostic Boundary}

STOP is defined as the highest-precedence diagnostic outcome. When STOP
conditions are met, the model explicitly forbids further diagnosis.

This boundary is not a heuristic choice; it is a formal consequence of
the executable specification. Continuing analysis beyond STOP would
require introducing assumptions not present in the ETS.

\subsection{Implications for Diagnostic Integrity}

By making structural failure modes explicit, $K_{\mathrm{EA}}$ prevents
the silent continuation of analysis beyond its domain of validity.

This does not imply that recovery is impossible, nor does it prescribe
any course of action. It states only that, within the model, diagnosis
is no longer defined once STOP conditions are met.

The distinction between diagnosable degradation and non-diagnosable
structural failure is central to the integrity of the enterprise
continuum formulation.
